\documentclass[11pt,a4paper]{article}
\usepackage[utf8]{inputenc}
\usepackage{amsmath,amssymb,amsthm}
\usepackage{listings}
\usepackage{xcolor}
\usepackage[margin=2.5cm]{geometry}
\usepackage{hyperref}

\lstset{
  language=C++,
  basicstyle=\ttfamily\small,
  keywordstyle=\color{blue}\bfseries,
  commentstyle=\color{green!60!black},
  stringstyle=\color{red!70!black},
  numbers=left,
  numberstyle=\tiny\color{gray},
  breaklines=true,
  frame=single,
  tabsize=4,
  showstringspaces=false
}

\title{IOI 2000 -- Car Parking}
\author{Solution}
\date{}

\begin{document}
\maketitle

\section{Problem Statement}

A parking lot has $N$ spaces. Cars are initially arranged in some configuration
(a permutation with one empty space, denoted 0), and must be rearranged to a target
configuration. Each move drives one car into the (unique) empty space.
Find the minimum number of moves.

\section{Solution Approach}

\subsection{Permutation Cycle Decomposition}

The initial and target configurations define a permutation $\pi$: for each space $i$,
the item currently at $i$ needs to move to some target position. Decompose $\pi$ into
cycles.

\begin{itemize}
  \item A \textbf{fixed point} (cycle of length 1) costs 0 moves.
  \item A cycle of length $k$ containing the empty space costs $k - 1$ moves.
        The empty space can rotate the cycle in $k - 1$ steps.
  \item A cycle of length $k$ \emph{not} containing the empty space costs $k + 1$ moves:
        1 move to bring the empty space into the cycle, then $k$ moves to rotate all
        $k$ elements into place (the empty space ends up back outside the cycle).
\end{itemize}

\begin{proof}[Proof of the $k + 1$ bound]
To rotate a cycle of length $k$ that does not contain the empty space, we must first
move the empty space to some position in the cycle (1 move), then perform $k$ moves
to shift each element one position along the cycle. Total: $k + 1$. This is optimal
because each of the $k$ misplaced elements must move at least once, and one additional
move is needed to bring the empty space into the cycle.
\end{proof}

\subsection{Algorithm}

\begin{enumerate}
  \item Build the permutation from initial to target configuration.
  \item Find all cycles using standard cycle-finding.
  \item For each non-trivial cycle ($k \ge 2$): add $k - 1$ if it contains the empty
        space, or $k + 1$ otherwise.
\end{enumerate}

\section{C++ Solution}

\begin{lstlisting}
#include <cstdio>
#include <vector>
#include <map>
using namespace std;

int main() {
    int N;
    scanf("%d", &N);

    vector<int> initial(N + 1), target(N + 1);
    int emptyPos = -1;

    for (int i = 1; i <= N; i++) {
        scanf("%d", &initial[i]);
        if (initial[i] == 0) emptyPos = i;
    }
    for (int i = 1; i <= N; i++)
        scanf("%d", &target[i]);

    // Map each car to its target position
    map<int, int> targetPos;
    int emptyTarget = -1;
    for (int i = 1; i <= N; i++) {
        if (target[i] == 0) emptyTarget = i;
        else targetPos[target[i]] = i;
    }

    // Build permutation: perm[i] = where the item at position i should go
    vector<int> perm(N + 1);
    for (int i = 1; i <= N; i++) {
        if (initial[i] == 0)
            perm[i] = emptyTarget;
        else
            perm[i] = targetPos[initial[i]];
    }

    // Find cycles and compute total moves
    vector<bool> visited(N + 1, false);
    int totalMoves = 0;

    for (int i = 1; i <= N; i++) {
        if (visited[i]) continue;
        int len = 0;
        bool containsEmpty = false;
        int cur = i;
        while (!visited[cur]) {
            visited[cur] = true;
            if (cur == emptyPos) containsEmpty = true;
            len++;
            cur = perm[cur];
        }
        if (len <= 1) continue; // fixed point
        if (containsEmpty)
            totalMoves += len - 1;
        else
            totalMoves += len + 1;
    }

    printf("%d\n", totalMoves);
    return 0;
}
\end{lstlisting}

\section{Complexity Analysis}

\begin{itemize}
  \item \textbf{Time complexity:} $O(N)$. Each position is visited exactly once during
        cycle detection.
  \item \textbf{Space complexity:} $O(N)$ for the permutation and visited arrays.
\end{itemize}

\end{document}
